% Section 2.3, from lectures/L02/appendix/b_exercises.md.

\section{Exercises}
\label{c2:sec:exercises}\label{c2:app:b}
These exercises reinforce the concepts from \secref{c2:app:a}.

\begin{aside}
\textbf{How to check your work.} Write each exercise set as its own application under
\file{lectures/L02/exercises}, with the file names given there, and run \code{make test} from the
repository root. A set’s tests switch on as soon as its files exist, and a set you have not started
is reported as \code{SKIP}. Each application grows through its set, and the tests check where it
ends up. Every set is built twice: once as you wrote it, and once with an extra overlay from the
test suite merged on top, which moves pins, turns polarities around, and adds or removes nodes. A
program that takes its hardware from the devicetree prints the second expected output without a
change; a program with a pin number or a polarity written into its C code does not.

The solutions are published in the companion repository, under
\file{lectures/L02/appendix/solutions}, a few days after the lecture. Try each exercise yourself
before you read them.
\end{aside}

Remember \code{source env.sh} in every new terminal. Build and run an exercise from the repository
root:

\begin{shell}
west build -p -b native_sim/native/64 lectures/L02/exercises/exercise1
west build -t run
\end{shell}

\noindent To build with the same extra overlay as the tests, from the repository root:

\begin{shell}
west build -p -b native_sim/native/64 lectures/L02/exercises/exercise1 -- \
    -DEXTRA_DTC_OVERLAY_FILE=$PWD/lectures/L02/exercises/test/exercise1-extra.overlay
\end{shell}

\noindent Print with \code{printk()} in every set, so that the output has no timestamps other than
the ones the exercise asks for.

% ---------------------------------------------------------------------
\exerciseset{The board’s LED}
\label{c2:set:1}
Create an application named \code{exercise1} to work in for this exercise set, with a
\code{CMakeLists.txt} and \code{src/main.c} as in Chapter~\ref{c1:ch}, and a \code{prj.conf} that
enables GPIO (\code{CONFIG_GPIO=y}).

\exercise{Finding the LED}{Code}
\label{c2:ex:1.1}
In this exercise you will get the board’s LED from the devicetree.

\task{Tasks}
\begin{itemize}
  \item In \code{main.c}, get a \code{struct gpio_dt_spec} for the \code{gpios} property of the node with
    the alias \code{led0}, at file scope.
  \item In \code{main()}, check that the LED’s controller is ready. If it is not, print \code{GPIO
    controller <name> is not ready} and return.
  \item Print the LED’s pin, its controller’s name, and whether it is active high or active low, from the
    spec’s flags:
\begin{console}
LED: pin 0 on gpio_emul, active high
\end{console}
  \item Configure the pin as an output, initially inactive.
\end{itemize}

\task{Reflection}
\begin{itemize}
  \item Find \code{led0} in \code{deps/zephyr/boards/native/native_sim/native_sim.dts}. Which line is the
    alias, and which is the label? Which of the two did your code use?
  \item Where does the name \code{gpio_emul} come from?
  \item What would the same program print on a board whose \code{led0} is declared
\code{<&gpio1 5 GPIO_ACTIVE_LOW>}, and what in your code would change?
\end{itemize}

\exercise{Blinking as the devicetree says}{Code}
\label{c2:ex:1.2}

\task{Tasks}
\begin{itemize}
  \item Create \code{app.overlay} in \code{exercise1}, adding a \code{zephyr,user} node under the root with
    two properties:
\begin{itemize}
  \item \code{blink-ms = <200 300>;}: how long the LED is on, and then off, in milliseconds.
  \item \code{blinks = <3>;}: how many times to blink.
\end{itemize}
  \item In \code{main.c}, read \code{blink-ms} into a \code{static const int} array at file scope, and use
    \code{blinks} directly where you need it.
  \item After configuring the pin, blink the LED \code{blinks} times. For each blink:
\begin{itemize}
  \item Turn the LED on with \code{gpio_pin_set_dt()}, print \code{LED on, pin <level> at <uptime> ms}, and
    sleep for the on time.
  \item Turn the LED off, print \code{LED off, pin <level> at <uptime> ms}, and sleep for the off time.
\end{itemize}
\code{<level>} is \code{high} or \code{low}: what the emulated controller is driving on the pin,
from \code{gpio_emul_output_get()} in \code{<zephyr/drivers/gpio/gpio_emul.h>}.
  \item Finally print \code{Done after <blinks> blinks}.
\end{itemize}

The output must be:

\begin{console}
LED: pin 0 on gpio_emul, active high
LED on, pin high at 0 ms
LED off, pin low at 210 ms
LED on, pin high at 520 ms
LED off, pin low at 730 ms
LED on, pin high at 1040 ms
LED off, pin low at 1250 ms
Done after 3 blinks
\end{console}

\noindent The tests build the application a second time with
\file{lectures/L02/exercises/test/exercise1-extra.overlay}, which points \code{led0} at a new,
active-low LED on pin 5, and blinks it twice, for 50 and 150 ms. The output must then be:

\begin{console}
LED: pin 5 on gpio_emul, active low
LED on, pin low at 0 ms
LED off, pin high at 60 ms
LED on, pin low at 220 ms
LED off, pin high at 280 ms
Done after 2 blinks
\end{console}

\noindent \textbf{Hint:} \code{DT_PATH(zephyr_user)} is the node identifier of \code{/zephyr,user}.
An array property is an initializer: \code{static const int blink_ms[] = DT_PROP(node, blink_ms);}.

\task{Reflection}
\begin{itemize}
  \item You wrote no binding for \code{zephyr,user}. How did the build know the types of its properties?
  \item Remove \code{blinks} from the overlay and build. Explain the error. Then make the program blink once
    when \code{blinks} is missing, without adding it back.
  \item In the second build, the LED is \emph{on} while the pin is \emph{low}. Which line of your code
    decided that, and why did it not need to change?
  \item In the second build, find \code{blink-ms} in \code{build/zephyr/zephyr.dts}. How does the file say
    which overlay the value came from, and why is it not \code{app.overlay}, which also sets it?
\end{itemize}

% ---------------------------------------------------------------------
\exerciseset{A button of your own}
\label{c2:set:2}
Create an application named \code{exercise2} to work in for this exercise set, with the same
\code{CMakeLists.txt} and \code{prj.conf} as \code{exercise1}.

\exercise{Describing the hardware}{Code}
\label{c2:ex:2.1}
In this exercise you will add hardware the board does not have, with an overlay.

\task{Tasks}
Create \code{app.overlay} in \code{exercise2}, which adds:

\begin{itemize}
  \item A second LED, as a child of the board’s existing \code{leds} node: pin 2 of \code{gpio0}, active
    high, with the node label \code{led1}.
  \item A node named \code{buttons} under the root, with \code{compatible = "gpio-keys"}, holding one
    button: pin 1 of \code{gpio0}, active low with a pull-up, with the node label \code{button0}.
  \item The aliases \code{led1} for the new LED and \code{sw0} for the button.
\end{itemize}

Use exactly these node labels: the test suite’s overlay refers to them.

\task{Reflection}
\begin{itemize}
  \item Why does the new LED need no \code{compatible} property of its own?
  \item Build the application and open \code{build/zephyr/zephyr.dts}. What are the button’s flags, in hex,
    and how are they made up?
\end{itemize}

\exercise{The callback}{Code}
\label{c2:ex:2.2}

\task{Tasks}
In \code{main.c}:

\begin{itemize}
  \item Get \code{gpio_dt_spec}s for the aliases \code{led1} and \code{sw0}, and check that both are ready.
  \item Print \code{Button on pin <pin>, LED 1 on pin <pin>}.
  \item Configure LED 1 as an output, initially inactive, and the button as an input.
  \item Configure an interrupt on the button for when it becomes \textbf{active}, and add a callback for it.
  \item In the callback, count the press, toggle a \code{bool} that says whether LED 1 is on, set LED 1 with
    \code{gpio_pin_set_dt()}, and print:
\begin{console}
Press <count> at <uptime> ms, LED 1 <on|off>, pin <high|low>
\end{console}
where the pin’s level is read back with \code{gpio_emul_output_get()}.
\end{itemize}

\exercise{The finger}{Code}
\label{c2:ex:2.3}

\task{Tasks}
Nobody can press an emulated button, so \code{main()} does it:

\begin{itemize}
  \item Write a function \code{static void press(bool pressed)} that sets the button’s input with
    \code{gpio_emul_input_set()}. It takes whether the button is pressed, and must work out the
    \textbf{level on the pin} from the button’s flags.
  \item After setting up the callback, \code{main()} presses the button three times. Before each press it
    sleeps for the next value of \code{static const int gap_ms[] = \{100, 200, 200\};}, then presses,
    sleeps 50 ms, and releases.
  \item Finally print \code{<count> presses counted}.
\end{itemize}

The output must be:

\begin{console}
Button on pin 1, LED 1 on pin 2
Press 1 at 110 ms, LED 1 on, pin high
Press 2 at 380 ms, LED 1 off, pin low
Press 3 at 650 ms, LED 1 on, pin high
3 presses counted
\end{console}

\noindent The second build merges \file{lectures/L02/exercises/test/exercise2-extra.overlay}, which
moves the button to pin 7 as an active-high button with a pull-down, and LED 1 to pin 3 as an
active-low LED. The output must then be:

\begin{console}
Button on pin 7, LED 1 on pin 3
Press 1 at 110 ms, LED 1 on, pin low
Press 2 at 380 ms, LED 1 off, pin high
Press 3 at 650 ms, LED 1 on, pin low
3 presses counted
\end{console}

\noindent \textbf{Hint:} An active-low button is pressed when its pin is low.

\task{Reflection}
\begin{itemize}
  \item Make \code{press()} always drive the pin low for a press, and run the second build. At what times
    are the presses counted now, and why?
  \item Print \code{k_is_in_isr()} in the callback. What does it print, and what would the same callback run
    in on a microcontroller? What does that forbid in the callback?
  \item Nothing happens when the finger releases the button. Why not?
  \item The button’s input level is high before the first press, although nothing in your code set it. What
    did?
  \item Build with \code{-DDTC_OVERLAY_FILE=} and the path to \code{exercise2-extra.overlay}, instead of
    \code{-DEXTRA_DTC_OVERLAY_FILE=}. What happens, and why?
\end{itemize}

% ---------------------------------------------------------------------
\exerciseset{A binding of your own}
\label{c2:set:3}
Create an application named \code{exercise3} to work in for this exercise set, with the same
\code{CMakeLists.txt} and \code{prj.conf} as \code{exercise1}.

\exercise{The binding}{Code}
\label{c2:ex:3.1}
In this exercise you will define a kind of node of your own.

\task{Tasks}
Create \code{dts/bindings/qacademy,alarm.yaml} in \code{exercise3}: a binding for the compatible
\code{qacademy,alarm}, a temperature alarm with an LED, which includes \code{base.yaml} and has
these properties:

\begin{booktable}{@{}lllL@{}}
  \thead{Property} & \thead{Type} & \thead{Required} & \thead{Meaning} \\
  \midrule
  \code{alarm-name} & string & yes & Human-readable name. \\
  \code{threshold} & int & yes & The temperature, in degrees Celsius, at which the alarm goes on. \\
  \code{hysteresis} & int & no, default 2 & How many degrees below the threshold the temperature must fall before the alarm goes off. \\
  \code{severity} & string & yes & One of \code{info}, \code{warning} or \code{critical}, in that order. \\
  \code{led-gpios} & phandle-array & yes & The LED that shows the alarm. \\
\end{booktable}

Create \code{dts/bindings/vendor-prefixes.txt}, registering the vendor prefix \code{qacademy} (the
prefix, a tab, and a vendor name).

\exercise{The alarms}{Code}
\label{c2:ex:3.2}

\task{Tasks}
Create \code{app.overlay} in \code{exercise3}, adding three nodes under the root, in this order.
Every alarm's LED is on \code{gpio0} and active high, and \code{alarm-fan} is the only node with
\code{status = "disabled"}:

\begin{narrowtable}{@{}lllllll@{}}
  \thead{Node name} & \thead{Label} & \thead{\code{alarm-name}} & \thead{\code{threshold}} &
    \thead{\code{hysteresis}} & \thead{\code{severity}} & \thead{LED pin} \\
  \midrule
  \code{alarm-overheat} & \code{overheat} & \code{Overheat} & 70 & not set & \code{critical} & 10 \\
  \code{alarm-warm} & \code{warm} & \code{Warm} & 45 & 5 & \code{warning} & 11 \\
  \code{alarm-fan} & \code{fan} & \code{Fan} & 25 & not set & \code{info} & 12 \\
\end{narrowtable}

\task{Reflection}
\begin{itemize}
  \item Remove \code{severity} from one node and build. Then set it to \code{"fatal"} and build. Which step
    of the build reports each error, and has any C been compiled yet?
  \item Delete \code{vendor-prefixes.txt} and build. What changes?
\end{itemize}

\exercise{The table}{Code}
\label{c2:ex:3.3}

\task{Tasks}
In \code{main.c}:

\begin{itemize}
  \item Define a \code{struct alarm} holding the name, threshold, hysteresis, severity (as an \code{int}), a
    \code{struct gpio_dt_spec} for the LED, and a \code{bool} saying whether the alarm is on.
  \item Build an array of them, \code{alarms}, with \textbf{one entry for each enabled \code{qacademy,alarm}
    node}, using \code{DT_FOREACH_STATUS_OKAY()}. Store the severity as its index in the binding’s
    \code{enum}, and keep a table of the names \code{"info"}, \code{"warning"} and \code{"critical"} to
    print it.
  \item Print \code{<count> alarms}, and then, for each alarm:
\begin{console}
<name>: <severity> at <threshold> C, hysteresis <hysteresis>, LED pin <pin>
\end{console}
and configure its LED as an output, initially inactive.
  \item Go through the readings \code{static const int readings[] = \{20, 50, 42, 75, 66, 30\};}. For each,
    print \code{Reading <t> C}, and then for each alarm, in table order:
\begin{itemize}
  \item An alarm that is off goes on if the reading is at or above its threshold.
  \item An alarm that is on goes off if the reading is below its threshold minus its hysteresis.
  \item When an alarm changes, set its LED, and print, indented by two spaces:
\begin{console}
  <name> (<severity>) <on|off>, pin <pin> <high|low>
\end{console}
\end{itemize}
\end{itemize}

The output must be:

\begin{console}
2 alarms
Overheat: critical at 70 C, hysteresis 2, LED pin 10
Warm: warning at 45 C, hysteresis 5, LED pin 11
Reading 20 C
Reading 50 C
  Warm (warning) on, pin 11 high
Reading 42 C
Reading 75 C
  Overheat (critical) on, pin 10 high
Reading 66 C
  Overheat (critical) off, pin 10 low
Reading 30 C
  Warm (warning) off, pin 11 low
\end{console}

\noindent The second build merges \file{lectures/L02/exercises/test/exercise3-extra.overlay}, which
deletes the \code{warm} node, enables \code{fan}, and changes \code{overheat}’s threshold to 60 and
hysteresis to 10. The output must then be:

\begin{console}
2 alarms
Overheat: critical at 60 C, hysteresis 10, LED pin 10
Fan: info at 25 C, hysteresis 2, LED pin 12
Reading 20 C
Reading 50 C
  Fan (info) on, pin 12 high
Reading 42 C
Reading 75 C
  Overheat (critical) on, pin 10 high
Reading 66 C
Reading 30 C
  Overheat (critical) off, pin 10 low
\end{console}

\noindent \textbf{Hint:} \code{DT_ENUM_IDX(node_id, severity)} is the index of the node’s
\code{severity} in the binding’s \code{enum}. The compatible, written as a C token, is
\code{qacademy_alarm}.

\task{Reflection}
\begin{itemize}
  \item Overheat’s hysteresis is 2, although its node does not set it. Where does the 2 come from, and what
    would happen without it in the binding?
  \item In what order are the alarms in the table? Find the list in
    \code{build/zephyr/include/generated/zephyr/devicetree_generated.h}. Can a program rely on that
    order?
  \item \code{alarm-fan} is disabled in the first build. Which instance number does it have, and does
    \code{DT_NODE_EXISTS(DT_NODELABEL(fan))} give 1 or 0?
  \item Add \code{BUILD_ASSERT(DT_PROP(DT_NODELABEL(warm), hysteresis) == 5);} to \code{main.c}, and run the
    second build. What happens, and why is the error not an assertion failure?
  \item The table is not \code{const}, because each alarm’s \code{active} changes. Which fields could be
    \code{const}, and how would you split the struct so that they are?
\end{itemize}

% ---------------------------------------------------------------------
\exerciseset{Reading the devicetree}
\label{c2:set:4}
No code in this set: answer the questions by looking. Build \code{exercise1} and \code{exercise2}
into separate build directories first, e.g. with \code{-d build/exercise1} and \code{-d
build/exercise2}.

\exercise[fillaccent2]{The merged tree}{Reflection}
\label{c2:ex:4.1}

\task{Tasks}
Answer, using \code{zephyr.dts} in the \code{exercise2} build:

\begin{itemize}
  \item Which file and line defined the alias \code{sw0}? And the alias \code{led0}?
  \item What does \code{/chosen} say the console is, and what is the node’s compatible?
\end{itemize}

\exercise[fillaccent2]{The generated header}{Reflection}
\label{c2:ex:4.2}

\task{Tasks}
Answer, using \code{devicetree_generated.h}:

\begin{itemize}
  \item What is the node identifier of the button in \code{exercise2}, and what does \code{DT_ALIAS(sw0)}
    expand to, step by step?
  \item What is the ordinal of \code{/gpio_emul} in the \code{exercise1} build, and in the \code{exercise2}
    build? Why can they differ?
  \item Which ordinals does the button’s node depend on directly, and why those?
\end{itemize}

\exercise[fillaccent2]{Two errors, one number}{Reflection}
\label{c2:ex:4.3}

\task{Tasks}
\begin{itemize}
  \item Build \code{exercise1} with \code{-DEXTRA_DTC_OVERLAY_FILE=} and the path to
    \code{lectures/L02/examples/blinky/broken.overlay}. Which error do you get, and which node does its
    number name?
  \item Build \code{exercise1} with \code{-DCONFIG_GPIO=n} instead. Which error do you get now, and what is
    the difference in what went wrong?
\end{itemize}
